% =============================================================================
% STANDALONE RESEARCH PAPER - Energy Consumption Pattern Analysis
% Self-contained LaTeX file - Just compile and get PDF!
% =============================================================================

\documentclass[11pt,a4paper]{article}

% Packages
\usepackage[utf8]{inputenc}
\usepackage[T1]{fontenc}
\usepackage{amsmath,amssymb}
\usepackage{graphicx}
\usepackage{booktabs}
\usepackage[hyphens]{url}
\usepackage[colorlinks=true,citecolor=blue,linkcolor=blue,urlcolor=blue]{hyperref}
\usepackage{geometry}
\geometry{margin=1in}

% Simple bibliography (no biber needed)
\usepackage{natbib}
\bibliographystyle{plain}

% Graphics path - Images in same folder
\graphicspath{{./}}

% Title
\title{\textbf{Shape-First Behavioral Segmentation of Household Energy Consumption: A Reproducible Framework}}

\author{
  Shantanu \and Aarna Srivastava \and Johnson Victor Yalangi \and Harsh Rathi \and Varun Srivastava\\
  [Institution]
}

\date{September 2026}

% =============================================================================
\begin{document}
\maketitle

% =============================================================================
% ABSTRACT
% =============================================================================
\begin{abstract}
Household energy consumption segmentation typically groups consumers by total usage magnitude, overlooking temporal behavioral patterns. We present a reproducible framework separating consumption magnitude from temporal behavior through 51 scale-invariant features derived from normalized 24-hour load profiles. The framework applies PCA (95\% variance, 10 components) and evidence-based K-selection combining silhouette, Calinski-Harabasz, Davies-Bouldin, and multi-seed stability. Validated on 200 consumers over 365 days with known ground-truth archetypes, the framework achieves K=4 with ARI=0.813 against hidden labels. Multi-seed stability yields ARI=0.995, temporal stability ARI=0.882. Ablation study shows magnitude-only features achieve near-zero behavioral recovery (ARI≈0.00) despite highest silhouette (0.52), confirming magnitude-timing separation. All code and data publicly available.

\textbf{Keywords:} household energy, clustering, PCA, K-means, behavioral segmentation, reproducible research
\end{abstract}

% =============================================================================
% 1. INTRODUCTION
% =============================================================================
\section{Introduction}

Smart meter deployment enables high-resolution monitoring of household electricity consumption. Traditional clustering groups households by \emph{how much} energy they use (magnitude), but demand response requires understanding \emph{when} they use energy (timing). Two households with identical daily load shapes but different scales should cluster together; households with similar totals but opposite peak hours should not.

\subsection{Research Gap}

Existing studies face challenges: (1) raw features allow magnitude to dominate; (2) K chosen heuristically; (3) validation lacks ground truth; (4) stability unreported; (5) limited reproducibility.

\subsection{Contributions}

\begin{enumerate}
\item 51 scale-invariant behavioral features separating magnitude from timing
\item Evidence-based K selection with multi-criterion composite rule
\item Controlled validation: ARI=0.813 on synthetic data with hidden archetypes
\item Comprehensive stability: multi-seed (ARI=0.995), temporal (ARI=0.882)
\item Rigorous ablation: magnitude-only features fail (ARI≈0.00)
\item Complete reproducibility infrastructure
\end{enumerate}

% =============================================================================
% 2. METHODOLOGY
% =============================================================================
\section{Methodology}

\subsection{Dataset}

Synthetic dataset with 200 consumers, 365 days, 4 behavioral archetypes (daytime, evening, flat, weekend). Magnitude drawn independently of archetype enabling validation of magnitude-timing separation. Ground truth hidden during all modeling.

\subsection{Feature Engineering}

For each consumer $i$, normalized shape: $\tilde{S}_{i,h} = \mu_{i,h} / \sum_{h'} \mu_{i,h'}$ where $\mu_{i,h}$ is mean hourly consumption. This makes shape invariant to scalar multiplication.

51 behavioral features:
\begin{itemize}
\item 24 normalized hourly bins
\item 11 timing features (period shares, peak timing, weekend ratio)
\item 9 shape descriptors (entropy, Gini, harmonics, wavelets)
\item 4 variability metrics (peak-to-avg, CV, skewness, kurtosis)
\item 3 dispersion metrics
\end{itemize}

\subsection{PCA and Clustering}

Standardize features, apply PCA with 95\% variance threshold → 10 components retained (95.05\% variance). K-means with K=2-10 candidates. Select K via composite rule: normalize silhouette, Calinski-Harabasz, inverted Davies-Bouldin; filter by cluster balance (≥5\%) and stability (ARI≥0.60); take smallest K within 5\% tolerance of best composite score.

% =============================================================================
% 3. RESULTS
% =============================================================================
\section{Results}

\subsection{K Selection}

K=4 selected with composite score 0.944. Silhouette=0.328, Calinski-Harabasz=96.6, Davies-Bouldin=1.169.

\begin{figure}[h]
\centering
\includegraphics[width=0.8\textwidth]{k_selection_metrics.png}
\caption{K-Selection Metrics: (a) Silhouette, (b) Calinski-Harabasz, (c) Davies-Bouldin, (d) Elbow. Selected K=4.}
\end{figure}

\subsection{Ground Truth Recovery}

ARI=0.813, NMI=0.828 against hidden archetypes. Recovery peaks at K=4, exactly where evidence-based rule selected.

\begin{figure}[h]
\centering
\includegraphics[width=0.6\textwidth]{archetype_recovery.png}
\caption{Ground Truth Recovery: ARI and NMI vs K. Peak at K=4 matches selection.}
\end{figure}

\begin{figure}[h]
\centering
\includegraphics[width=0.7\textwidth]{archetype_crosstab.png}
\caption{Confusion Matrix at K=4. Three archetypes recover near-perfectly (evening 47/50, flat 50/50, weekend 49/50).}
\end{figure}

\subsection{Stability Analysis}

\textbf{Multi-seed:} Mean pairwise ARI=0.9947±0.0071 across 10 random restarts. Near-perfect reproducibility.

\textbf{Temporal:} Mean ARI=0.882 across 4 quarterly segments. Behavioral groups persist despite 47\% seasonal magnitude variation.

\begin{figure}[h]
\centering
\includegraphics[width=0.7\textwidth]{longitudinal_cluster_stability.png}
\caption{Temporal Stability: Quarterly segment ARI vs full window. Mean=0.882.}
\end{figure}

\subsection{Ablation Study}

Five feature sets compared (scale, shape, summary, behavioral, combined). Scale-only achieves highest silhouette (0.521) but zero behavioral recovery (ARI=-0.004). Behavioral features achieve ARI=0.614.

\begin{figure}[h]
\centering
\includegraphics[width=0.8\textwidth]{ablation_comparison.png}
\caption{Ablation Study: Scale features fail behavioral recovery despite high internal quality.}
\end{figure}

\textbf{Critical Finding:} Internal metrics alone cannot validate behavioral segmentation. High silhouette on magnitude simply confirms consumers with different totals are well-separated—the wrong research question.

\subsection{Robustness Across 20 Datasets}

Behavioral features: mean ARI=0.641±0.115. Statistical test: behavioral vs scale p=1.9×10$^{-6}$ (highly significant).

\begin{figure}[h]
\centering
\includegraphics[width=0.7\textwidth]{seed_robustness.png}
\caption{Seed Robustness: Behavioral features achieve consistent recovery across 20 independent datasets.}
\end{figure}

\subsection{PCA and Explainability}

10 components explain 95.05\% variance. Top loadings:
\begin{itemize}
\item PC1: profile\_ramp +0.92, peak\_concentration +0.88
\item PC2: night\_day\_ratio +0.92, afternoon\_share -0.87
\item PC3: morning pattern (hour 7-8 shape)
\end{itemize}

\begin{figure}[h]
\centering
\includegraphics[width=0.75\textwidth]{explained_variance.png}
\caption{PCA: Explained and cumulative variance. 10 components retained at 95\% threshold.}
\end{figure}

SHAP analysis: clusters distinguished by semantically appropriate features (weekend cluster driven by weekend\_ratio 0.250 importance).

\begin{figure}[h]
\centering
\includegraphics[width=0.75\textwidth]{shap_cluster_importance.png}
\caption{SHAP Feature Importance per Cluster. CV surrogate accuracy=0.985.}
\end{figure}

% =============================================================================
% 4. DISCUSSION
% =============================================================================
\section{Discussion}

\subsection{Principal Findings}

\textbf{1. Magnitude-timing separation is critical.} Ablation proves magnitude-only features fail behavioral segmentation (ARI≈0) despite high silhouette (0.52).

\textbf{2. Evidence-based K aligns with ground truth.} Composite rule selects K=4 using only internal metrics, independently validated by ARI peak at K=4.

\textbf{3. Robust across time and seeds.} Stability (0.995) and temporal persistence (0.882) demonstrate reproducible behavioral groupings.

\textbf{4. Internal metrics insufficient.} Scale-only arm proves silhouette alone cannot validate behavioral meaning.

\subsection{Practical Implications}

Evening-peaking clusters (54.5\%) can receive time-of-use pricing. Flat profiles (26\%) require baseline reduction. Weekend-heavy (31\%) respond to weekend-specific rates. Targeted programs vs uniform enrollment.

\subsection{Limitations}

Synthetic primary dataset with simplified archetypes. Real-world pathway implemented but not executed. Moderate scale (200 consumers). K-means spherical assumptions. PCA linearity. Single flagship run without bootstrap CIs.

% =============================================================================
% 5. CONCLUSION
% =============================================================================
\section{Conclusion}

We demonstrate reproducible behavioral segmentation of household energy consumption by separating magnitude from temporal patterns through 51 scale-invariant features, evidence-based K selection, and comprehensive stability validation. Ablation conclusively shows magnitude-only features fail (ARI≈0) while behavioral features achieve ARI=0.813 on controlled synthetic data. Stability across seeds (0.995) and time (0.882) confirms robust behavioral types.

Contributions: (1) shape normalization methodology, (2) evidence-based K selection, (3) controlled validation, (4) comprehensive stability, (5) rigorous ablation, (6) complete reproducibility. Real-world validation remains essential future work.

\subsection{Reproducibility}

Repository: \url{https://github.com/shaxntanu/Energy-Consumption-Pattern-Analysis-using-PCA-and-K-Means}

Configuration hash: \texttt{99c7a6631340d301} (flagship 365-day run)

All code, data generation, and artifacts publicly available.

% =============================================================================
% REFERENCES (Simplified)
% =============================================================================
\section*{References}

\begin{thebibliography}{9}

\bibitem{kwac2014} Kwac et al. (2014). Household Energy Consumption Segmentation Using Hourly Data. IEEE Trans. Smart Grid.

\bibitem{lazar2021} Lazar et al. (2021). Investigating Underlying Drivers of Variability in Residential Energy Usage Patterns. arXiv:2102.11027.

\bibitem{afzalan2020} Afzalan et al. (2020). Two-Stage Clustering of Household Electricity Load Shapes. arXiv:2008.04293.

\bibitem{jolliffe2002} Jolliffe (2002). Principal Component Analysis. Springer.

\bibitem{lloyd1982} Lloyd (1982). Least Squares Quantization in PCM. IEEE Trans. Information Theory.

\bibitem{arthur2007} Arthur \& Vassilvitskii (2007). K-means++: The Advantages of Careful Seeding. ACM-SIAM SODA.

\bibitem{rousseeuw1987} Rousseeuw (1987). Silhouettes: A Graphical Aid to Cluster Analysis. J. Comp. Applied Math.

\bibitem{calinski1974} Caliński \& Harabasz (1974). A Dendrite Method for Cluster Analysis. Comm. Statistics.

\bibitem{peng2011} Peng (2011). Reproducible Research in Computational Science. Science.

\end{thebibliography}

% =============================================================================
\end{document}
